\subsection{Uniform tails for every original flag length}
\label{global-gram-continuation-length-tail}

The nonidentity terms with zero Gram determinant require a tail
estimate for every flag length and its inverse.  The first-length
and simultaneous-ratio estimates supply this assertion with the
original range of moments.  We first explain why a given square
integrable Bessel coefficient has the required normalization.

\begin{lemma}\label{global-gram-finite-bessel-realization}
Let $\sigma$ be an irreducible unitary representation on $H$, let
$S$ be a closed unimodular subgroup of a unimodular Lie group $G$,
and let $\rho$ be a finite-dimensional unitary representation of $S$.
Use the given Haar measures and their quotient measure.  Suppose
$B:H^\infty\to W$ is a continuous smooth $S$-equivariant map.
If, for some $v_0\in H^\infty$,
\[
       0<\int_{S\backslash G}\|B\sigma(g)v_0\|^2dq(g)<\infty,
\]
then $v\mapsto B\sigma(\cdot)v$ extends to a positive scalar
multiple of an isometric $G$-embedding of $H$ into the original
$L^2(S\backslash G,W)$.
\end{lemma}
\begin{proof}
Let $D$ consist of the smooth vectors whose displayed coefficients
belong to $L^2$, and define $T$ on $D$ by that coefficient.
The domain is invariant under $G$ and contains the span of the
translates of $v_0$.  It is therefore dense in $H$.

We prove closability before using Schur's lemma.  For a compactly
supported smooth section $\phi$ of the quotient bundle, choose
local sections of $G\to S\backslash G$, a finite partition of
unity over its support, and a compactly supported smooth function
on $S$ of Haar integral one.  The Weil formula then writes
$\langle Tv,\phi\rangle$ as a finite sum of expressions
\[
                 \xi\bigl(B\sigma(\eta)v\bigr),
        \qquad \eta\in C_c^\infty(G),\quad \xi\in W^*.
\]
Each expression is bounded in the Hilbert norm of $v$.
Indeed compact smooth convolution maps $H$ continuously into
$H^\infty$: for any word $D_1$ in the Lie algebra, its derivative
is convolution by the corresponding derivative of $\eta$, with
operator norm at most that derivative's $L^1$ norm.  Continuity
of $B$ uses only finitely many such seminorms.
If $v_j\to0$ in $H$ and $Tv_j\to F$ in $L^2$, all these
test pairings vanish.  Density of the compact smooth sections
gives $F=0$.

The closure $\overline T$ still intertwines the unitary actions.
Every spectral projection of $\overline T^*\overline T$ therefore
commutes with $\sigma$ and is zero or the identity.  The projections
of the intervals $[0,N]$ increase strongly to the identity, so one
already equals the identity.  Schur's lemma now gives
$\overline T^*\overline T=aI$ for a finite $a>0$.
Thus the closure is defined on all of $H$.  The same compact-test
identity identifies its value at every smooth $v$ with
$B\sigma(\cdot)v$ as a distribution, hence as the original smooth
coefficient.  Square integrability of derivatives is a consequence,
not a premise, of this argument.
\end{proof}

Let $r\ge0$, $k\ge1$, $n=2r+k+1$, $m=2r+2$, and let
$G=U(k+r,r+1)$ in the original unitary Witt coordinates.
Write $e_0,\ldots,e_{r-1}$ for its standard isotropic flag,
$f_0,\ldots,f_{r-1}$ for the opposite Witt vectors with
$h(e_i,f_j)=\delta_{ij}$, and let $z$ lie in the central
space, with Hermitian norm $-T<0$.  Put $e_r=z$ and
\[
 L_i(g)=\operatorname{dist}\left(g^{-1}e_i,
                    \operatorname{span}_{j<i}g^{-1}e_j\right)
                       \quad(0\le i\le r).
\]
The distance is Euclidean.  Retain the original flag Bessel
subgroup, its compact target, character, and quotient measure.
Let $\Phi(g)=B\sigma(g)w$, where $w$ is a unit vector in the
lowest compact type and $0<\|\Phi\|_2^2<\infty$.

Here are the mathematical size hypotheses used below.  They are
clauses or consequences of the original standing hypotheses.
Let $E/F$ be the original CM extension and $f=[F:\mathbb Q]$.
Write $b=e^L$, $L>1$, $d=\operatorname{rd}(E)$, and $g=g_1g_2$.
Assume $C_g,C_T\ge0$, $C_\ell\ge1000(1+C_1)+10^8$,
$C_1>0$, and
\begin{gather*}
 m^4\le f\le k,\qquad C_\ell\le m,\qquad
 1\le d/2,\qquad \log d\le L/C_\ell,\\
 1\le g\le b^{2C_g},\qquad 1\le T\le b^{C_Tm},\\
 |\lambda_a|\le b^2,\qquad |u_i|\le b^U,
                    \qquad U=\max(2,C_u).
\end{gather*}
The Harish--Chandra parameter $\lambda$ satisfies
$\lambda_a-(n-1)/2\in\mathbb Z$.  Its first $k-1$ signs are
positive and its remaining $m$ signs alternate $+,-$.
It interlaces the compact parameter $u$ on its first $k-1$
coordinates.  For a
number $B_v$ the first $k-1$ coordinates of $\lambda$ belong
to $[0,B_v/8]$, and the remaining $m$ coordinates belong to
$[-2B_v,-B_v]$, where $10^8n^7\le B_v\le b^2$.
Every adjacent gap is larger than $C_\ell k m^4+n m^{10}$,
and every noncompact root gap is at least $C_\ell n^{19}f^5L$.
The local character is exactly
\[
 \psi(n_0)=\exp\left(i\operatorname{Re}
                 \left(\frac{2\pi}{(d/2)g}\,\ell(n_0)\right)\right),
\]
where $\ell$ is the sum of the original adjacent flag coordinates,
including its final coordinate involving $z$.
These definitions retain the original scalar $d/2$.
\cite{global-gram-length-tail-source}

\begin{theorem}\label{global-gram-original-all-length-tail}
Put
\[
 A=U+5,\qquad A_c=2C_g+2,\qquad
 C_H=C_T+8A+A_c.
\]
Let $q\ge1$ be an integer such that
\[
 q\le C_1kr,\qquad mq\le Q,\qquad
                   mq\le\frac{B_v}{10^8n^4},
\]
where $Q$ is the original real shift allowance.  Only the integer
shifts of size at most $2q\le Q$ are used.  For every $t>0$,
\begin{equation}\label{global-gram-original-length-tail-bound}
 \int_{\{g:\,\exists i,\ \max(L_i(g),L_i(g)^{-1})
                                  \ge t b^{C_Hm}\}}
       \|\Phi(g)\|^2dq(g)
       \le m t^{-2q}\|\Phi\|_2^2.
\end{equation}
In particular this gives the original assertion with constant
$2m$.  For $q=0$ that assertion follows directly from positivity.
The constants $A,A_c,C_H$ are independent of $C_1$, the field,
the word, the compact unit vector, and every Haar scaling.
\end{theorem}
\begin{proof}
First use the standard character
$\exp(2i\operatorname{Im}\ell)$ and the vector
$z=\sqrt T\,e_-$.  Set $\beta=b^A$ and
$K=n^2T\beta^4(6B_v)^{2m}\le T\beta^{6m}$.
We check the reserves for the established first-length and ratio
estimates, rather than increasing them.
Since $m\ge2C_1$ and $k\ge m^4$,
\[
 q\le C_1kr\le km^2/4\le k^2\le n^2,
 \qquad
 2rq\le C_1km^2/2\le km^3/4\le k^2\le n^3.
\]
The moment of the first length used below has order $2q$ and
therefore uses exactly $2q$ integer shifts and level $4q\le4n^2$.
The simultaneous ratio moment has total order $2iq\le2rq\le n^3$.
Also $2q\le mq\le Q$ and
$2q\le B_v/(10^8n^4)\le B_v/100$.
The displayed condition $mq\le B_v/(10^8n^4)$ is a separate
original hypothesis.  It also follows in the present range from
$mq\le n^3$ and $B_v\ge10^8n^7$; no factor $n^4$ is inferred
from a bound involving only $mk$.
The adjacent-gap bound gives gaps at least $100(2q+1)$
before shifting.  After every shift used here they are at least two,
the two parameter intervals lie in their original $B_v/100$
enlargements, and the noncompact gap remains at least
$8n^{19}$.  Thus the original signed Pieri and normalized Bessel
families exist throughout these shifts, by the already proved
family construction.  No additional tensor reserve is requested.
The derivative-energy bound applies with base $\beta$:
the lowest compact Casimir radius is at most $4nB_v$ and
$\beta\ge3B_v$.  All quotient measures are the given ones.
\cite{current-construction-record}

We also retain the character sign.  The rank-two Lie algebra
element in the ratio calculation has adjacent coordinate
$itL_{i+1}^2$.  The present character differentiates to
$+2iL_{i+1}^2$, whereas the earlier ratio display used the
negative of this scalar.  Taking squared norms gives exactly
the same ratio inequality.  In the first-length calculation
the signed tensor transfer and projection estimates do not
depend on that sign.  The adjacent root factors in its Capelli
evaluation have the same moduli $1$ and $\sqrt T$, so the
degree-normalized lowering bound remains $T\beta^4$.
This is the direct sign check used in the quantitative-seed
proof; no representation or compact parameter is conjugated.

Choose a compact-highest unit vector $w_0$ in the lowest compact
type.  This vector has level zero.  Write
$F_v=\sqrt{d_\sigma}\,\mathcal B\sigma(\cdot)v$ for the
subgroup-period normalized realization, and put
$R_i=L_{i+1}/L_i$.  The established first-length and mixed-ratio
estimates give
\[
 \int L_0^{4q}\|F_{w_0}\|^2dq\le K^{2q},\qquad
 \int\prod_{j<i}R_j^{4q}\|F_{w_0}\|^2dq
                                  \le\beta^{8iq}.
\]
Their zero-order norm is one when $\mathcal B\ne0$.
Since $L_i=L_0\prod_{j<i}R_j$, Cauchy--Schwarz gives
\begin{equation}\label{global-gram-length-positive-moment}
 \int L_i^{2q}\|F_{w_0}\|^2dq
                  \le K^q\beta^{4iq}
                  \le b^{(C_T+8A)mq}.
\end{equation}
The residual last vector has Hermitian norm $-T$, so
$L_r\ge\sqrt T$.  Consequently
$L_i^{-1}=L_r^{-1}\prod_{j=i}^{r-1}R_j$ and the ratio
moment of total order $(r-i)q$ gives
\begin{equation}\label{global-gram-length-negative-moment}
 \int L_i^{-2q}\|F_{w_0}\|^2dq
       \le T^{-q}\beta^{4(r-i)q}\le b^{2Amq}.
\end{equation}
This includes $i=r$ with an empty product.

Neither estimate requires all lowest-type vectors to have level
zero.  The weights $L_i^{2q}$ and $L_i^{-2q}$ are invariant
under right multiplication by the maximal compact group.
The compact orbit of $w_0$ spans the lowest compact type.
The weighted integrals are finite on this entire finite-dimensional
space, by the triangle inequality and their finiteness on that
orbit.  Each is a compact-invariant Hermitian form, hence a
scalar multiple of the original form.  Thus the same bounds
hold for every unit vector in that type.

For the given Bessel map, Lemma~\ref{global-gram-finite-bessel-realization}
gives a bounded embedding.  It legitimately intertwines convolution
by the integrable lowest-vector matrix-coefficient projector.
The original projector kernel is
$F_w(g)F_w(h)^*$, as proved in the quotient-mass argument.
Its range is therefore $\mathbb C F_w$; if $\mathcal B=0$
its range is zero.  The given positive-norm coefficient is fixed
by this projector, so the zero case is excluded and its unit
normalization differs from $F_w$ by a unit scalar.
Equations~\eqref{global-gram-length-positive-moment} and
\eqref{global-gram-length-negative-moment} therefore hold relative
to its same unweighted norm.  This argument does not assume
boundedness of a smooth Bessel functional on $H$.

Finally transport the original character and vector exactly.
A central unitary isometry sends $\sqrt T e_-$ to $z$ and
fixes the isotropic flag.  Next set
\[
 c=\frac{g d}{2\pi},\qquad
 \eta_i=(-ic)^{r-i},\qquad
 D_\eta e_i=\eta_i e_i,
 \quad D_\eta f_i=\overline{\eta_i}^{-1}f_i,
 \quad D_\eta z=z.
\]
Conjugation multiplies each adjacent root coordinate by $-ic$,
so the original character becomes the standard one.
Push the original subgroup and quotient measures by these
changes of variables; the Weil formula introduces no unspecified
factor.  The compact target and the input vector are unchanged.
For the $i$th original length the change is exactly
$L_i\mapsto c^{-(r-i)}L_i$.
Since
\[
 \max(c,c^{-1})\le b^{A_c},\qquad 2r<m,
\]
each of the two moment bounds costs at most an additional
$b^{A_cmq}$.  Notice that this normalized local-coordinate
scale contains the root discriminant.

Each positive and inverse moment is now at most
$b^{C_Hmq}\|\Phi\|_2^2$.
Markov's inequality at the threshold $t b^{C_Hm}$ bounds
each corresponding tail by $t^{-2q}\|\Phi\|_2^2$.
There are $2(r+1)=m$ tails.  Their union gives
\eqref{global-gram-original-length-tail-bound}.
\end{proof}

The original shift allowance is
$Q=\min\{c_q h_{\rm NC}/\sqrt m,4\Delta\}/C_q$.
Choose $C_q=10^8$ and $c_q=C_q^{-1}$ once, before choosing
$C_1$.  Thus the two branches are exactly
$h_{\rm NC}/(C_q^2\sqrt m)$ and $4\Delta/C_q$.
The preceding theorem retains both branches and the original
condition $mq\le Q$.  The original interval condition says
$C_qn^7\le B_v$, $0\le v_l\le B_v/8$, and
$B_v\le b_j\le2B_v$.  The actual ordered parameter is
\[
             \lambda=(v_1,\ldots,v_{k-1},-b_1,\ldots,-b_m).
\]
It gives exactly the two parameter intervals used above, and
$B_v\le b_1\le b^2$.  The remaining size assumptions follow
from the original standing data after increasing the final fixed
threshold $C_\ell$, before constructing the arithmetic data.
In particular $|u_i|\le b^U$ follows from compact interlacing,
$u_1\le b^{C_u}$ and $u_k>-b_1\ge-b^2$.
The negative adjacent gaps and the cloud-to-negative gap give
the displayed noncompact gap, using
$b_1\ge C_\ell mk\Delta$ and $\Delta\ge(kfL)^{C_\ell}$.
Indeed $n\le2k$, $k\ge2$, $f\ge1$, $L>1$ and
$C_\ell\ge40$ give
$mk(kfL)^{C_\ell}\ge2^{19}k^{19}f^5L\ge n^{19}f^5L$.
The constants $c_q,C_q,C_H$ depend only on the fixed original
size constants, whereas the final threshold may also depend on
$C_1$.  This is the original order of quantifiers.
The proof is independent of the permitted Haar factor $\theta$.
When $C_1=0$ only $q=0$ is allowed, and when $C_1<0$ there is
no allowed nonnegative $q$ in the actual range $r\ge1$.
It supplies the original all-length tail assertion, including its
quantifier order.  The local bound for the remaining zero-Gram
orbital operators and their global sum are still separate obligations.
